% ============================================================
%  Appendix A — CII Regulatory Framework
% ============================================================
\section{CII Regulatory Framework}
\label{app:cii}

\subsection{Regulatory Background}
\label{app:cii:background}

The Carbon Intensity Indicator (\CII{}) was introduced by the International
Maritime Organization (\IMO{}) under MARPOL Annex VI, entering into force on
1 November 2022 and applying from 1 January 2023.  It forms part of the
short-term measures in the IMO Initial GHG Strategy and the revised 2023 GHG
Strategy targeting net-zero emissions by 2050.

Under \CII{}, all ships of 5{,}000 GT and above engaged in international
voyages must:
\begin{enumerate}
  \item Calculate their \emph{attained \CII{}} annually from reported fuel
    consumption and distance data;
  \item Compare it to the \emph{required \CII{}} derived from a
    ship-type-specific reference line;
  \item Receive a rating (A--E) based on their position relative to rating
    boundary multipliers;
  \item Implement a corrective action plan if rated D or E.
\end{enumerate}

\subsection{CII Formula by Ship Type}
\label{app:cii:formula}

The attained \CII{} is computed as:
\begin{equation}
  \ciiatt = \frac{\sum_j CF_j \cdot FC_j}{\text{Capacity} \times D_T},
  \label{eq:cii_general}
\end{equation}
where $CF_j$ is the CO$_2$ conversion factor for fuel type $j$, $FC_j$ is
the mass of fuel $j$ consumed (tonnes), and $D_T$ is the total distance
sailed (nautical miles).  The \emph{Capacity} denominator is
ship-type-specific:

\begin{table}[h]
\centering
\caption{Capacity parameter by ship type for CII calculation
         (MEPC.352(78) / MEPC.354(78))}
\label{tab:cii_capacity}
\begin{tabular}{ll}
  \toprule
  Ship type & Capacity \\
  \midrule
  Bulk carrier & DWT \\
  Tanker        & DWT \\
  Container ship & DWT \\
  General cargo  & DWT \\
  \textbf{Ro-Ro cargo (vehicle carrier)} & \textbf{GT} \\
  \textbf{Ro-Ro cargo (non-vehicle carrier)} & \textbf{GT} \\
  \textbf{Ro-Ro passenger} & \textbf{GT} \\
  Cruise passenger ship & GT \\
  LNG carrier  & DWT \\
  \bottomrule
\end{tabular}
\end{table}

For \RoRo{} cargo vessels (non-vehicle carrier), the operative formula is:
\begin{equation}
  \ciiatt = \frac{FC \times CF}{f_i \times f_m \times \GT \times D_T},
  \label{eq:cii_roro}
\end{equation}
where $f_i$ and $f_m$ are the ice-class and capacity correction factors
respectively (\cref{tab:cii_corrections}).

\begin{table}[h]
\centering
\caption{Correction factors for RoRo CII (MEPC.352(78))}
\label{tab:cii_corrections}
\begin{tabular}{lll}
  \toprule
  Factor & Definition & Value for M/V Obbola \\
  \midrule
  $f_i$ & Ice-class factor (1.0 for non-ice class) & 1.0 \\
  $f_m$ & Capacity correction (1.05 for RoRo non-vehicle) & 1.05 \\
  \bottomrule
\end{tabular}
\end{table}

\subsection{Required CII and Rating Boundaries}
\label{app:cii:required}

The required \CII{} for year $Y$ is given by:
\begin{equation}
  \ciireq(Y) = \ciireq(2019) \times (1 - r_Y),
  \label{eq:cii_required}
\end{equation}
where $r_Y$ is the cumulative reduction factor for year $Y$ relative to
the 2019 reference.  For \RoRo{} non-vehicle carriers, the required CII
reference value is computed from the ship-type-specific regression:
\begin{equation}
  \ciireq(2019) = a \times C^{-c},
  \label{eq:cii_ref}
\end{equation}
with parameters $a$ and $c$ from MEPC.352(78) Table 1 and $C$ being the
capacity (\GT{} for \RoRo{}).

The five rating bands are defined by four boundary multipliers
$d_1 < d_2 < 1 < d_3 < d_4$ applied to the required \CII{}:

\begin{table}[h]
\centering
\caption{Rating boundary multipliers for RoRo cargo ships (non-vehicle
         carrier), from MEPC.339(76)}
\label{tab:rating_boundaries}
\begin{tabular}{ccc}
  \toprule
  Rating & Condition & Multiplier \\
  \midrule
  A & $\ciiatt < d_1 \times \ciireq$ & $d_1 = 0.76$ \\
  B & $d_1 \leq \ciiatt < d_2 \times \ciireq$ & $d_2 = 0.89$ \\
  C & $d_2 \leq \ciiatt < d_3 \times \ciireq$ & $d_3 = 1.08$ \\
  D & $d_3 \leq \ciiatt < d_4 \times \ciireq$ & $d_4 = 1.27$ \\
  E & $\ciiatt \geq d_4 \times \ciireq$ & --- \\
  \bottomrule
\end{tabular}
\end{table}

Note that the multipliers in \cref{tab:rating_boundaries} are
\RoRo{}-specific.  Generic vessel types use different multipliers
(e.g., $d_1 = 0.82$, $d_3 = 1.06$ for bulk carriers); applying generic
values to a \RoRo{} vessel would produce incorrect ratings.

\subsection{Confirmed and Projected Reduction Factors}
\label{app:cii:reduction}

\begin{table}[h]
\centering
\caption{CII reduction factors --- confirmed and projected}
\label{tab:reduction_factors}
\begin{tabular}{ccl}
  \toprule
  Year & Reduction (vs.\ 2019) & Status \\
  \midrule
  2023 & 5\,\%  & Confirmed (MEPC.338(76)) \\
  2024 & 7\,\%  & Confirmed \\
  2025 & 9\,\%  & Confirmed \\
  2026 & 11\,\% & Confirmed \\
  2027--2030 & TBD & Pending MEPC 83/84 \\
  \bottomrule
\end{tabular}
\end{table}

For scenarios S4 and S4b, this thesis assumes a 31\,\% reduction factor
for 2030, consistent with the IMO 2023 GHG Strategy's indicative 40\,\%
fleet-wide ambition and extrapolating the confirmed annual trajectory.
This is a scenario assumption and not a confirmed regulatory requirement.

\subsection{EU MRV Reporting and Data Collection}
\label{app:cii:mrv}

The EU Monitoring, Reporting and Verification (MRV) Regulation
(EU~2015/757, amended by EU~2023/957) requires ships above 5{,}000 GT
calling at EU ports to report per-voyage fuel consumption, CO$_2$
emissions, and distance sailed.  The reported data are submitted to the
THETIS-MRV platform and published annually.

For M/V Obbola, the EU MRV observation period covers December 2021
through August 2022, comprising 153 voyage legs.  The total reported
CO$_2$ emission is 16{,}342.18~t (15{,}727.27~t HFO + 614.91~t MGO),
yielding an actual 2022 attained \CII{} of 11.430253~\gGTnm{}
(Rating~B).
